\section{Related Work}\label{sec:related}

This section reviews relevant literature from four intersecting areas: spiking neural networks, predictive coding theory, attention mechanisms, and neuromorphic computation.

\subsection{Spiking Neural Networks and Markovian Dynamics}

Since Hodgkin and Huxley proposed the nonlinear ordinary differential equations (ODEs) that describe action potentials in the squid giant axon~\cite{hodgkin1952}, computational neuroscience has been built upon models grounded in physical and chemical principles. The high computational cost of the Hodgkin--Huxley formulation has made it difficult to apply in large-scale network simulations. As a practical alternative, the Leaky Integrate-and-Fire (LIF) model and its variants, such as the Izhikevich model~\cite{izhikevich2003} and the adaptive exponential integrate-and-fire model~\cite{brette2005}, have been widely adopted as the elementary units of spiking neural networks.

Although these models capture the bioenergetic efficiency and subthreshold dynamics of neuronal membranes~\cite{gerstner2014}, they are constrained by strictly Markovian dynamics: the state update at time $t$ depends only on the state at $t-1$ and the instantaneous input~\cite{abbott1999}. This memory-free character limits their ability to extract contextual information over extended temporal windows. Later work has introduced temporal hysteresis through short-term synaptic plasticity or slow recovery variables associated with spike-frequency adaptation~\cite{benda2003}, but these additions rely on local capacitive decay and cannot actively extract higher-order temporal dependencies in the way that modern recurrent neural networks or Transformers can.

Recent advances in directly-trained spiking neural networks have shown that SNNs can be trained end-to-end on large-scale tasks using surrogate gradients~\cite{eshraghian2021,neftci2019,fang2021}. The incorporation of learnable membrane time constants~\cite{fang2021} and differentiable spike functions~\cite{li2021differentiable} has brought SNN performance closer to that of artificial neural networks on image classification and event-stream tasks. St\"ockl and Maass~\cite{stockl2021} demonstrated that optimised single spiking neurons can classify images with high accuracy using just two spikes, highlighting the untapped potential of single-neuron computation. However, these improvements are achieved at the \textit{learning algorithm} level and do not modify the underlying \textit{single-neuron dynamics}: the neuron itself remains a Markovian integrator, and the temporal modelling burden is shifted entirely to the training procedure. TAN takes a different approach---it modifies the neuron's intrinsic dynamics to incorporate a non-Markovian temporal receptive field, while keeping the parameter set small (5 scalars) and requiring no training.

\subsection{Predictive Coding and the Free-Energy Principle}

How the brain processes temporal sequences has long been a central question in cognitive neuroscience. Friston's free-energy principle proposes that biological systems strive to minimize the surprise (prediction error) between internal generative models and incoming sensory signals---effectively reducing informational entropy~\cite{friston2010}. Within this framework, the cortex can be viewed as a hierarchical network that transmits prediction errors upward and predictions downward~\cite{keller2018,rao1999}. The formalisation of \textit{active inference}~\cite{friston2017} extends this idea into a unified principle of perception, action, and learning, in which biological agents minimise variational free energy over both internal beliefs and external interventions.

The free-energy framework connects naturally to the notion of surprise as a driver of neural response. Biologically, surprise signals are consistent with theories proposing that neural circuits continuously compare incoming sensory evidence against internal predictions and preferentially transmit prediction errors~\cite{srinivasan1982,friston2010}. The sparse coding hypothesis~\cite{olshausen1996} further argues that neural systems allocate their limited metabolic resources preferentially to ``unexpected entropic events''~\cite{attwell2001}, a principle that TAN implements through its surprise-driven gating, although TAN's deviation signal is a simplified local-window statistic rather than a full predictive-coding architecture.

\subsection{Attention Mechanisms in Deep Learning and SNNs}

In deep learning, sequence modeling has undergone a notable shift in recent years. The Transformer architecture, introduced by Vaswani et al.~\cite{vaswani2017}, employs a global self-attention mechanism via query-key-value interactions and has markedly influenced the paradigm of sequence processing~\cite{lin2022survey,tay2022efficient}. A number of recent studies have explored the integration of attention-like operations into spiking neural networks, for example by converting temporal synaptic weights into dynamic attention matrices~\cite{yao2021} or by introducing local self-attention pooling into deep SNNs~\cite{zheng2020,zhou2023spikformer,yao2023spikedriven}. The Spikformer architecture~\cite{zhou2023spikformer} explicitly replaces softmax self-attention with a spiking-based attention mechanism, achieving competitive performance on image classification. Spike-driven Transformers~\cite{yao2023spikedriven} further demonstrate the feasibility of fully spike-based attentional processing.

However, almost all of these efforts apply attention at the \textit{network layer} scale, using it primarily as a performance-boosting engineering tool across spatial features. They do not explore the temporal counterpart of the attention operation at the \textit{single-neuron} scale. In contrast, the TAN model proposed here does not simply attach a Transformer block to a macroscopic network. Instead, it combines a QKV-like temporal weighting mechanism with a deviation-driven gating signal at the single-neuron level. This perspective may offer a biologically interpretable dynamical mapping for self-attention mechanisms, although we do not claim that TAN is mathematically equivalent to a Transformer.

\subsection{Neuromorphic Computation and Artificial Life}

A natural motivation for spiking neural networks is their compatibility with neuromorphic hardware, which processes information via sparse, event-driven spikes rather than dense matrix multiplications. Intel's Loihi~\cite{davies2018} demonstrated on-chip learning and real-time inference at microwatt power budgets; its successor Loihi 2 (announced 2021) extended these capabilities with an open-source neuromorphic computing framework. BrainScaleS-2~\cite{pehle2022} and SpiNNaker 2 further extended the hardware landscape toward mixed-signal and digital neuromorphic platforms. Recent surveys~\cite{frenkel2023} note that the \textit{native computational primitive} of such hardware---a leaky integrator with threshold-based spike generation---provides the same building blocks that TAN extends with attention and gating.

A central challenge is that most current SNN learning algorithms require either (a) conversion from trained artificial networks, which sacrifices the online-adaptive advantage of SNNs, or (b) backpropagation-based surrogate gradient training, which demands substantial off-chip computation~\cite{eshraghian2021,neftci2019}. Backpropagation-free local learning rules such as STDP~\cite{bi1998} are biologically plausible but have so far lagged in task performance. TAN offers a complementary route: rather than improving the \textit{learning algorithm}, it modifies the \textit{native dynamics} of each neuron so that useful temporal behaviour emerges from the model itself, without requiring any training. This positions TAN as a candidate computational primitive for future neuromorphic systems that must adapt to noisy, non-stationary environments~\cite{roy2019}. However, a net energy advantage cannot be established without hardware-level energy measurements; TAN is spike-sparse but introduces additional internal computation (window storage, QKV operations, attention, gating).

The field of artificial life (ALife)~\cite{langton1995} aims to replicate emergent life-like behaviors through bottom-up computational approaches. Early thought experiments, such as Braitenberg's minimalist vehicles~\cite{braitenberg1984}, demonstrated that even simple sensor--motor cross-connections can give rise to behavioral patterns that evoke psychological categories such as fear, aggression, or phototaxis. Highly realistic simulations such as OpenWorm~\cite{szigeti2014} and the Blue Brain Project~\cite{markram2015} encounter a ``curse of complexity'': even when the topological wiring of the nervous system is carefully reproduced, the absence of precise data on dynamic synaptic weights often makes it challenging for the system to spontaneously generate robust, high-level life-like behaviors~\cite{gjorgjieva2016}. As demonstrated by systems such as the Game of Life~\cite{gardner1970}, complex emergent phenomena do not necessarily require individually complex components. Motivated by this line of thought, the present study returns to a minimalist mathematical mapping philosophy, employing the temporal attention dynamics of TAN rather than intricate ion channels.

% ============================================================
% 3. The Temporal Attention Neuron Model
% ============================================================
